式~\eqref{eq:subspace_hessian}--\eqref{eq:elliptic_group_condition} 在数学层面建立了"分组即子空间选择"的映射，但这一映射是否在实际优化中产生可观测的后果，需要从最简单的可精确验证情形开始。F1 是对角椭球型函数，其 Hessian 呈严格对角形，因此每个组的条件数恰好等于组内最大与最小曲率尺度之比，无需任何采样近似。表~\ref{tab:f1_condition} 汇总了在目标函数完全固定的前提下，仅改变分组所产生的条件数与终点差异。为紧凑起见，condition 列仅报告各诱导子问题条件数的均值。
